\chapter{Minimum Spanning Trees}
\chapterepigraph{To connect every node at least total cost, greedily add
the cheapest edge that joins two so-far-separate parts. The Disjoint Set
data structure makes ``are these already connected?'' almost free, and the
same structure -- run in reverse, or watched over time -- answers dynamic
connectivity questions too.}

\section{Greedy Connection with Disjoint Set}
\label{sec:cses-graph-mst-intro}

A minimum spanning tree (MST) connects all $n$ nodes with $n-1$ edges of
minimum total weight. Kruskal's algorithm sorts edges and adds each that
does not form a cycle, using Union-Find for the cycle test.

\begin{patternbox}[MST and Connectivity Problems]
\begin{itemize}
  \item \textbf{Cheapest way to connect everything} -- \textbf{Kruskal's
        MST} with Disjoint Set: Road Reparation (report impossible if the
        graph is disconnected).
  \item \textbf{Add edges over time, tracking components and the largest
        component} -- \textbf{incremental Disjoint Set}: Road
        Construction.
  \end{itemize}
\end{patternbox}

\begin{ideabox}[Union-Find Makes the Cycle Test O(1)]
Kruskal's correctness rests on always taking the cheapest edge between two
different components. Disjoint Set answers ``same component?'' in near
constant time (with path compression and union by size), and its size
counter directly gives the largest-component answer for Road Construction.
\end{ideabox}

\newpage

\section{Road Reparation}
\label{sec:cses-road-reparation}

\problemstatement{Given $n$ cities and $m$ possible roads, each with a
repair cost, choose roads of minimum total cost so that every city is
reachable from every other. Output the minimum cost, or \texttt{IMPOSSIBLE}
if the cities cannot all be connected.}

\begin{patternbox}
``Connect all cities at minimum total cost'' is a \textbf{minimum spanning
tree}. \textbf{Kruskal's algorithm} sorts roads by cost and adds each that
joins two different components, using Disjoint Set for the cycle test; if
fewer than $n-1$ roads are added, the graph is disconnected.
\end{patternbox}

\subsection*{Kruskal with Union-Find}

Sort all roads ascending by cost. Process each: if its endpoints are in
different components (checked via Disjoint Set), add it to the MST, union
the components, and accumulate the cost. Skip roads whose endpoints are
already connected (they would form a cycle). After processing, if exactly
$n-1$ roads were added, output the total; otherwise the graph could not be
fully connected, so print \texttt{IMPOSSIBLE}.

\begin{ideabox}[Cheapest Safe Edge First]
The MST is built greedily: the cheapest edge crossing between two
components is always safe to include. Disjoint Set makes ``same component?''
near-constant, and reaching $n-1$ edges certifies a spanning tree exists.
\end{ideabox}

\subsection*{Algorithm}

\begin{enumerate}
  \item Sort roads by cost. Initialise Disjoint Set.
  \item For each road, if it joins two components, union and add its cost;
        count added roads.
  \item If $n-1$ roads were added, output the total cost; else
        \texttt{IMPOSSIBLE}.
\end{enumerate}

\subsection*{Dry run}

\dryrun{$n=3$; roads $1\!-\!2(3)$, $2\!-\!3(5)$, $1\!-\!3(4)$.}

\begin{itemize}
  \item Sorted: $3,4,5$. Add $1\!-\!2(3)$; add $1\!-\!3(4)$; skip
        $2\!-\!3$ (cycle).
  \item Total $=3+4=7$, two roads $=n-1$: connected.
\end{itemize}

\subsection*{C++ implementation}

\begin{lstlisting}
#include <bits/stdc++.h>
using namespace std;

struct DSU {
    vector<int> parent, sz;
    DSU(int n) : parent(n + 1), sz(n + 1, 1) { iota(parent.begin(), parent.end(), 0); }
    int find(int x) { return parent[x] == x ? x : parent[x] = find(parent[x]); }
    bool unite(int a, int b) {
        a = find(a); b = find(b);
        if (a == b) return false;
        if (sz[a] < sz[b]) swap(a, b);
        parent[b] = a; sz[a] += sz[b];
        return true;
    }
};

int main() {
    int n, m;
    cin >> n >> m;
    vector<tuple<long long,int,int>> edges(m);
    for (auto& [w, a, b] : edges) cin >> a >> b >> w;
    sort(edges.begin(), edges.end());

    DSU dsu(n);
    long long total = 0;
    int used = 0;
    for (auto& [w, a, b] : edges)
        if (dsu.unite(a, b)) { total += w; used++; }

    if (used == n - 1) cout << total << "\n";
    else cout << "IMPOSSIBLE\n";
    return 0;
}
\end{lstlisting}

\begin{complexitybox}
\textbf{Time Complexity.} $O(m\log m)$ -- sorting dominates; unions are
near-constant. \textbf{Space Complexity.} $O(n+m)$.
\end{complexitybox}

\begin{takeawaybox}
Minimum-cost full connection is an MST; Kruskal plus Disjoint Set builds it
and, by counting added edges, detects disconnection. This is the canonical
network-design primitive.
\end{takeawaybox}

\section{Road Construction}
\label{sec:cses-road-construction}

\problemstatement{There are $n$ cities and initially no roads. Roads are
built one at a time ($m$ of them). After each new road, output the current
number of connected components and the size of the largest component.}

\begin{patternbox}
Adding edges and tracking connectivity over time is \textbf{incremental
Disjoint Set}. Maintain a running component count and the maximum component
size, both updated in $O(\alpha)$ whenever a union actually merges two
different components.
\end{patternbox}

\subsection*{Update counts on each real union}

Start with $n$ components, each of size $1$, and a maximum size of $1$.
For each new road, union its endpoints: if they were already in the same
component, nothing changes. If they were separate, the component count
drops by one, the merged component's size is the sum of the two, and the
maximum size is updated against it. Print the count and maximum after each
road.

\begin{ideabox}[Only Real Merges Change the Answer]
Union-Find's size counter gives the merged component's size immediately, so
the largest-component tracker updates in $O(1)$ per successful union. A
union that finds both endpoints already together leaves both statistics
untouched.
\end{ideabox}

\subsection*{Algorithm}

\begin{enumerate}
  \item Disjoint Set of $n$ singletons; $\textit{components}=n$,
        $\textit{maxSize}=1$.
  \item Per road: if the endpoints unite, decrement $\textit{components}$
        and update $\textit{maxSize}$ with the new size.
  \item Output $\textit{components}$ and $\textit{maxSize}$ after each road.
\end{enumerate}

\subsection*{Dry run}

\dryrun{$n=4$; roads $1\!-\!2$, $3\!-\!4$, $2\!-\!3$.}

\begin{itemize}
  \item After $1\!-\!2$: components $3$, max $2$.
  \item After $3\!-\!4$: components $2$, max $2$.
  \item After $2\!-\!3$: merges the two pairs -- components $1$, max $4$.
\end{itemize}

\subsection*{C++ implementation}

\begin{lstlisting}
#include <bits/stdc++.h>
using namespace std;

int main() {
    int n, m;
    cin >> n >> m;
    vector<int> parent(n + 1), sz(n + 1, 1);
    iota(parent.begin(), parent.end(), 0);
    function<int(int)> find = [&](int x) { return parent[x] == x ? x : parent[x] = find(parent[x]); };

    int components = n, maxSize = 1;
    while (m--) {
        int a, b; cin >> a >> b;
        a = find(a); b = find(b);
        if (a != b) {
            if (sz[a] < sz[b]) swap(a, b);
            parent[b] = a;
            sz[a] += sz[b];
            components--;
            maxSize = max(maxSize, sz[a]);
        }
        cout << components << " " << maxSize << "\n";
    }
    return 0;
}
\end{lstlisting}

\begin{complexitybox}
\textbf{Time Complexity.} $O((n+m)\,\alpha(n))$ -- near-linear.
\textbf{Space Complexity.} $O(n)$.
\end{complexitybox}

\begin{takeawaybox}
Online connectivity under edge additions is incremental Union-Find, with the
component count and largest-size maintained on each real merge. Disjoint Set
handles edge \emph{additions} gladly; edge \emph{deletions} are the hard
case needing offline or advanced structures.
\end{takeawaybox}

